\section{モジュール境界と最小手順}
\label{sec:modules}

\subsection{分割}

推奨するソース分割は次のとおりである。矢印は、記録側が作業側を呼ばない向きだけを許す。

\begin{verbatim}
state
  torsion params, admissible test, J, M
written_flow
  written_at, paper energy, momenta
working_models
  runaway, oscillator, sourced jets, quartic
layer_force
  gamma_s_equal, gravityLayerTerm, gravityLayerForce
channels
  pairChannelForce, coulombForce, pointForce
step_contract
  working_step, contract log
checks
  one file per claim kind
discrete_enum
  enumerate_admissible
\end{verbatim}

\texttt{step\_contract} は \texttt{written\_flow} と \texttt{layer\_force} と、必要なとき \texttt{channels} を呼ぶ。\texttt{written\_flow} は \texttt{layer\_force} を呼ばない。\texttt{checks} は本番の段から呼ばれない。

\subsection{最小手順}

実装の順序は次に固定する。後の段に進む前に、前の段のテストが通る状態にする。

\begin{enumerate}
\item \texttt{state} と許容判定。ランダムな六つ組を許容と非許容に分け、境界 \(\alpha+\beta=\pi/2\) を一点ずつ確認する。
\item \texttt{written\_at}。第\ref{sec:checks}章の内在角テスト。
\item 二体の \texttt{gravityLayerTerm} と反対称な \texttt{gravityLayerPush}。
\item \(N\) 体の \texttt{gravityLayerForce}。自己項が無いこと。
\item 契約ログ付きの \texttt{working\_step} を、粒子二つ、一軸の種一つで一段だけ進める。拒否の四条件をそれぞれ一度は踏む。
\item チャート比較の区間テストを、内在角ソルバとリンクせずに置く。
\item 余力があれば \texttt{enumerate\_admissible} を小さい \(N\) で実行し、点数が \(N^6\) 以下であることと、許容点が錐の中にあることを見る。
\end{enumerate}

\subsection{名前}

公開関数の名前は、Lean の識別子と揃える。作業模型だけ接頭辞 \texttt{working\_} を付ける。チャート専用のテストには接頭辞 \texttt{chart\_} を付ける。同じ綴りで別の式を返す別名を作らない。
